\documentclass{report}
\usepackage{amsmath,amssymb}
\setlength{\parindent}{0mm}
\setlength{\parskip}{1em}
\begin{document}
\begin{center}
\rule{6in}{1pt} \
{\large P. van Slingerland \\
{\bf A preconditioner for CG that does not need symmetry}}

Delft Institute of Applied Mathematics \\ Delft University of Technology \\ Mekelweg 4 \\ 2628 CD \\ Delft \\ The Netherlands
\\
{\tt P.vanSlingerland@tudelft.nl}\\
C. Vuik\end{center}

Although it is common to use a \emph{symmetric} preconditioner for the
conjugate gradient method, in this presentation, we demonstrate that a
\emph{non-symmetric} preconditioning strategy can be significantly more
efficient. The focus, motivation and conclusions of this research can be
summarized as follows.

This work is focused on an effective preconditioning strategy to increase
the efficiency of the Conjugate Gradient (CG) method for Symmetric and
Positive-Definite (SPD) linear systems resulting from Symmetric Interior
Penalty (discontinuous) Galerkin (SIPG) discretizations for diffusion
problems with extreme contrasts in the coefficients, such as those
encountered in oil reservoir simulations.

A Discontinuous Galerkin (DG) method can be thought of as a higher-order
generalization of the finite volume method that uses (discontinuous)
piecewise polynomials of degree $p$ rather than piecewise constants. As
such, it combines the best of both classical finite element methods and
finite volume methods, and it is particularly suitable for handling
non-matching grids and designing hp-refinement strategies. However, a
relevant drawback is that its resulting linear system is often
ill-conditioned and relatively large due to the large number of unknowns
per mesh element. In search of suitable iterative solution techniques,
much attention has been paid to subspace correction methods, such as
Schwarz domain decomposition; geometric \mbox{(h-)multigrid}; spectral
\mbox{(p-)multigrid}; and algebraic multigrid.

In particular, Dobrev et al. [\textit{Numer. Linear Algebra Appl}., 13
(2006), pp. 753--770] have proposed a spectral two-level preconditioner
that makes use of coarse corrections based on the solution approximation
with polynomial degree $p=0$. They have shown theoretically that this
preconditioner yields scalable convergence of the CG method (independent
of the mesh element diameter) for a large class of problems. Another nice
property is that the use of only two levels offers an appealing
simplicity. More importantly, the coefficient matrix that is used for the
coarse correction is quite similar to a matrix resulting from a central
difference discretization, for which very efficient solution techniques
are readily available.

However, two main issues remain when using this preconditioner for a SIPG matrix $A$:
first, \emph{two} smoothing steps must be applied during each iteration,
and the smoother needs to satisfy an inconvenient criterion to ensure
that the preconditioning operator is SPD. The second issue is that the
SIPG method involves a stabilizing penalty parameter, whose influence on
both $A$ and the preconditioner is not well understood for problems with
strongly varying coefficients. On the one hand, this parameter needs to
be chosen suffciently large to ensure that the SIPG method is stable and
convergent, and that $A$ is SPD. At the same time, it needs to be chosen
as small as possible to avoid an unnecessarily large condition number.
Known computable theoretical lower bounds [\textit{J. Comput. Appl.
Math.}, 206 (2007), pp. 843--872] are based on the ratio between the
\emph{global} maximum and minimum of the diffusion coefficient, and are
therefore impractical for our application.

To eliminate one of the two smoothing steps and the inconvenient
restriction on the smoother at the same time, we have cast the spectral
two-level preconditioner into the deflation framework by using the
analysis by Tang et al. [\textit{J. Sci. Comput.}, 39 (2009), pp.
340--370]. Additionally, we have studied the potential of a penalty
parameter that is based on \emph{local} values of the diffusion
coefficient, instead of the usual strategy to use one global constant for
the entire domain.

During the presentation, we discuss how and why the proposed spectral
two-level deflation method can be incorporated in a CG algorithm in the
form of an \emph{asymmetric} preconditioning operator. Furthermore, we
demonstrate numerically how the resulting iterative scheme performs for
diffusion problems with strongly varying coefficients, for both a
constant and a diffusion-dependent penalty parameter.

Our main findings can be summarized as follows: by reformulating the
spectral two-level preconditioner as a deflation method and using a
diffusion-dependent penalty parameter, the CG method can become over
$100$ times faster, while retaining scalable convergence (independent of
the mesh element diameter). Additionally, the SIPG convergence is
significantly better.
Altogether, the SIPG penalty parameter can best be chosen
diffusion-dependent, and the efficiency of a symmetric two-level
preconditioner can be significantly improved by switching to an
asymmetric deflation variant.


\end{document}
